\scanpage{81}
\mspage{74}
% archive sheet number: 80
un petit ensemble de gén. de $M$ (car $M$ est accessible), que $W$ soit stable par $\varinjlim$ filtrantes. Je devrais demander à Thomason \unsure{s'il|s'îl} \unsure{connait|connaît} un contreexemple.

c) Les deux exemples précédents étaient des catégories ordonnées. Comme exemple typique de cat. non ordonnée, où je ne vois pas ce qui se passe, je signale $I=B_{\unsure{G|C}}$, $G$ un groupe (\unsure{donc} un seul objet, avec un groupe d'autom. $G$). Aussi le cas de $\unsure{\mathbb{\Delta}|\Delta}[1]=\{\mathbb{\Delta}^0,\mathbb{\Delta}^1\}$ (sous-cat. pleine de $\mathbb{\Delta}$) n'est pas clair pour moi.

2) On aimerait pouvoir traiter de façon duale l'existence des $u_!^W$, à défaut de mieux. Dans l'hypothèse de départ du \S\ 4 sur $(M,W)$, cela consiste donc à changer \unsure{des $+$ et des $-$|des $+$ en $-$}. Cette fois, on devra supposer l'existence d'une quillénisation donc \uline{exacte à droite} pour $W$. Ce qui est plus exigeant, semble-t-il, qu'une telle structure exacte à gauche, dans le cas envisagé), avec
\[
\mathrm{TF}=F\cap C,
\]
